\chapter{URDF, Robot Modeling, and Coordinate Frames}
\label{ch:urdf}

\section{The source is URDF, not Xacro}
The runtime model is the plain XML file
\file{gp7_ros2/gp7_description/urdf/gp7.urdf}. Although the Gazebo package
manifest declares an Xacro dependency, no Xacro file is present and the task
launch does not invoke Xacro. Instead, Python parses and edits XML text before
passing the expanded description to Gazebo and robot-state publisher. This is
an important repository-specific distinction.

\section{Link and joint semantics}
A URDF \code{link} is a rigid body with optional visual geometry, collision
geometry, and inertial properties. A URDF \code{joint} connects one parent
link to one child link. Its \code{origin} is the pose of the joint frame
relative to the parent link; its \code{axis} is expressed in the joint frame.
For a revolute joint, the scalar joint position rotates the child-side chain
about that axis. Fixed joints preserve a transform but add no degree of
freedom. The \code{Kinematics} parser follows parent/child links from
\code{base_link} to \code{grasp_frame}; at each joint it applies the fixed
origin and then the revolute rotation.

The authoritative chain includes the six arm revolute joints, a fixed flange,
fixed tool0, fixed gripper base, one driver and five mimic finger joints, and a
fixed grasp frame. The extra \code{base} link is fixed to \code{base_link} at
z=\SI{0.33}{m}; it is not another actuated joint. The world-to-base transform
is identity.

\section{Visual, collision, and inertial models}
Each GP7 arm link references separate visual and collision STL files. The
visual files preserve appearance; collision geometry determines contact and
may be lower detail. The seven adaptive gripper DAE files are referenced for
visuals and, in this URDF, the same DAE path also appears in collision tags.
The working URDF uses scale 1.5 for the gripper meshes while leaving arm
meshes at scale 1; this is a project model adaptation. The environment includes
a 20 m square plane. Box A is a
static collision/visual box of size $(0.30,0.30,0.20)$ m. Box B is a dynamic
$(0.05)^3$ m cube with mass \SI{0.05}{kg}.

The URDF gives every modeled robot link a mass of \SI{1}{kg} and diagonal
inertia entries of $0.01$ with zero products of inertia. These are repeated
placeholder-like values, not validated CAD mass properties. They can permit
simulation but should not be presented as an identified physical inertia
model. Box B has diagonal inertia entries $2.08333\times10^{-5}$ kg m$^2$,
consistent with a uniform cube using its stated mass and dimensions. Box A is
static and does not need a dynamic inertial model for this task.

\section{TF tree and time-varying transforms}
\code{robot_state_publisher} consumes the URDF and joint-state messages and
publishes transforms for the link tree. Fixed joint transforms do not change;
the six arm and active gripper joint transforms depend on current joint
positions. In a conventional TF tree, each child has one parent, matching the
URDF chain. Gazebo's object model pose is not automatically the same stream as
the robot's TF transform: this task separately subscribes to the Gazebo
Transport pose vector for \code{box_b}.

The fixed identity transform from \code{world} to \code{base_link} aligns
the planning base with the simulation world. This is why the world-coordinate
Box A target can be passed to the URDF-based FK/IK without an additional
world/base calibration transform. That assumption would have to be changed if
the robot were spawned at a nonidentity pose.

\section{Actual zero-pose DH geometry}
The DH frames are deliberately not identified with all URDF link frames. At
zero joint angles, the documented frame origins in base coordinates are
\begin{align*}
 o_0&=(0,0,0), &o_1&=(0.040,0,0.330),\\
 o_2&=(0.040,0,0.775), &o_3&=(0.040,0,0.815),\\
 o_4&=o_5=o_6=(0.480,0,0.815)\quad\text{(m)}.
\end{align*}
Frame 3 lies at the intersection of the extended elbow axis and wrist axes,
not at the URDF \code{link_4_r} origin. The coincident wrist-frame origins
reflect intersecting wrist axes. The terminal DH frame 6 is still distinct
from the tool0/grasp frame.

\begin{figure}[htbp]
\centering\begin{tikzpicture}[scale=1.0,>=Stealth]
  \draw[->,gray,thick] (-.2,0)--(6.6,0) node[right]{base $+X$};
  \draw[->,gray,thick] (0,-.2)--(0,5.5) node[above]{base $+Z$};
  % Origins projected into X-Z; y=0 at the zero configuration.
  \draw[very thick,gray!55] (0,0)--(0,1.65)--(.2,1.65)--(.2,3.875)--(.2,4.075)--(2.4,4.075);
  \foreach \x/\z/\lab in {0/0/F0,.2/1.65/F1,.2/3.875/F2,.2/4.075/F3,2.4/4.075/{F4,F5,F6}} {
    \fill (\x,\z) circle (2pt);
    \node[anchor=west,font=\small] at (\x+.08,\z+.10) {\lab};
  }
  % Frame 0 and 1 axes projected in plane.
  \draw[->,red!75!black,thick] (0,0)--(.7,0) node[right,font=\scriptsize]{$x_0=+X$};
  \draw[->,blue!75!black,thick] (0,0)--(0,.55) node[above,font=\scriptsize]{$z_0=+Z$};
  \draw[->,red!75!black,thick] (.2,1.65)--(.85,1.65) node[right,font=\scriptsize]{$x_1=+X$};
  \draw ( .2,1.65) circle (0.10); \fill (.2,1.65) circle (1.5pt);
  \node[anchor=west,font=\scriptsize,blue!75!black] at (.32,1.86){$z_1=+Y$ out};
  % F2 has z into page; x2 is vertical.
  \draw[->,red!75!black,thick] (.2,3.875)--(.2,4.35) node[above,font=\scriptsize]{$x_2=+Z$};
  \draw (.2,3.875) circle (.10); \draw (.15,3.825)--(.25,3.925); \draw (.15,3.925)--(.25,3.825);
  \node[anchor=west,font=\scriptsize,blue!75!black] at (.32,3.70){$z_2=-Y$ into};
  % Frame 3 x3=+Z, z3=-X.
  \draw[->,red!75!black,thick] (.2,4.075)--(.2,4.55);
  \draw[->,blue!75!black,thick] (.2,4.075)--(-.38,4.075) node[left,font=\scriptsize]{$z_3=-X$};
  \node[anchor=west,font=\scriptsize] at (.32,4.52){$x_3=+Z$};
  % Wrist center: axes differ in 3D while origin coincides.
  \draw[->,red!75!black,thick] (2.4,4.075)--(2.4,4.55);
  \draw[->,blue!75!black,thick] (2.4,4.075)--(1.9,4.075);
  \node[anchor=west,font=\scriptsize] at (2.5,4.48){$x_{4,5,6}=+Z$};
  \node[anchor=north,font=\scriptsize,align=center] at (2.4,3.96){$z_4=-Y$; $z_5,z_6=-X$\\same wrist origin};
  \node[anchor=north west,font=\scriptsize,align=left] at (3.5,2.9){
    $F_0=(0,0,0)$\\
    $F_1=(.040,0,.330)$\\
    $F_2=(.040,0,.775)$\\
    $F_3=(.040,0,.815)$\\
    $F_{4,5,6}=(.480,0,.815)$\\[2mm]
    $y_i=z_i\times x_i$\\
    Circled dot: out of page\\
    Cross: into page\\[2mm]
    DH frame 6 is not the grasp frame.\\
    Use the separate $T_{6g}$ tool transform.
  };
\end{tikzpicture}

\caption{Standard-DH origin schematic projected into base X--Z. All pictured
origins lie at $y=0$ at home. Out-of-plane axis signs and origin coordinates
are annotated. The diagram is a kinematic construction, not a mesh rendering.}
\label{fig:dh-frames}
\end{figure}

\section{URDF fixed transforms near the tool}
Three fixed transforms matter in the DH-to-runtime comparison:
\begin{enumerate}
 \item \code{joint_6_t-flange}: $p=(0.080,0,0)$ m, zero RPY.
 \item \code{flange-tool0}: $p=(0,0,0)$ and RPY
 $(3.14159265,-1.57079632,0)$ rad.
 \item \code{tool0_to_grasp_frame}: $p=(-0.0009111,0.0123,0.14598)$ m,
 zero RPY.
\end{enumerate}
Replacing the decimal RPY angles by exact $\pi$ values is a convenient
derivation approximation but slightly changes the transform. The verification
script reads the exact URDF values.

\section{ros2\_control model embedded in URDF}
The URDF's \code{ros2_control} element declares a \code{GazeboSystem} with
\code{gz_ros2_control/GazeboSimSystem}. Each arm joint and the gripper
driver has a position command interface and position/velocity state
interfaces. The five mimic joints expose state interfaces but no command
interfaces. The external plugin loads the controller manager inside the
simulator and connects these abstract interfaces to Gazebo joint state and
command components. The project YAML chooses the actual controller instances.

\section{CAD exports versus runtime model}
The repository contains at least two CAD-exported GP7 URDFs in addition to the
runtime file: \file{GP7 (Solidworks 2020)/GP7 components/GP7/GP7.urdf} and
\file{GP7 (Solidworks 2020)/URDF/Build/Build.urdf}. They have distinct tree,
path, and mesh context and are not silently interchangeable. The task launch,
DH verification, and regression tests all target
\file{gp7_ros2/gp7_description/urdf/gp7.urdf}. A separate CAD-fidelity study
would need to compare origins, scaling, joint axes, and mesh coordinate frames
before claiming equivalence.
